%!TeX root=../emmatop.tex
\chapter[Chapter \thechapter]{}
\lettrine[lraise=0.3]{M}{r} Knightley was to dine with them—rather against the inclination of Mr Woodhouse, who did not like that any one should share with him in Isabella's first day. Emma's sense of right however had decided it; and besides the consideration of what was due to each brother, she had particular pleasure, from the circumstance of the late disagreement between Mr Knightley and herself, in procuring him the proper invitation.

She hoped they might now become friends again. She thought it was time to make up. Making-up indeed would not do. \textit{She} certainly had not been in the wrong, and \textit{he} would never own that he had. Concession must be out of the question; but it was time to appear to forget that they had ever quarrelled; and she hoped it might rather assist the restoration of friendship, that when he came into the room she had one of the children with her—the youngest, a nice little girl about eight months old, who was now making her first visit to Hartfield, and very happy to be danced about in her aunt's arms. It did assist; for though he began with grave looks and short questions, he was soon led on to talk of them all in the usual way, and to take the child out of her arms with all the unceremoniousness of perfect amity. Emma felt they were friends again; and the conviction giving her at first great satisfaction, and then a little sauciness, she could not help saying, as he was admiring the baby,

<What a comfort it is, that we think alike about our nephews and nieces. As to men and women, our opinions are sometimes very different; but with regard to these children, I observe we never disagree.>

<If you were as much guided by nature in your estimate of men and women, and as little under the power of fancy and whim in your dealings with them, as you are where these children are concerned, we might always think alike.>

<To be sure—our discordancies must always arise from my being in the wrong.>

<Yes,> said he, smiling—<and reason good. I was sixteen years old when you were born.>

<A material difference then,> she replied—<and no doubt you were much my superior in judgment at that period of our lives; but does not the lapse of one-and-twenty years bring our understandings a good deal nearer?>

<Yes—a good deal \textit{nearer}.>

<But still, not near enough to give me a chance of being right, if we think differently.>

<I have still the advantage of you by sixteen years' experience, and by not being a pretty young woman and a spoiled child. Come, my dear Emma, let us be friends, and say no more about it. Tell your aunt, little Emma, that she ought to set you a better example than to be renewing old grievances, and that if she were not wrong before, she is now.>

<That's true,> she cried—<very true. Little Emma, grow up a better woman than your aunt. Be infinitely cleverer and not half so conceited. Now, Mr Knightley, a word or two more, and I have done. As far as good intentions went, we were \textit{both} right, and I must say that no effects on my side of the argument have yet proved wrong. I only want to know that Mr Martin is not very, very bitterly disappointed.>

<A man cannot be more so,> was his short, full answer.

<Ah!—Indeed I am very sorry.—Come, shake hands with me.>

This had just taken place and with great cordiality, when John Knightley made his appearance, and <How d'ye do, George?> and <John, how are you?> succeeded in the true English style, burying under a calmness that seemed all but indifference, the real attachment which would have led either of them, if requisite, to do every thing for the good of the other.

The evening was quiet and conversable, as Mr Woodhouse declined cards entirely for the sake of comfortable talk with his dear Isabella, and the little party made two natural divisions; on one side he and his daughter; on the other the two Mr Knightleys; their subjects totally distinct, or very rarely mixing—and Emma only occasionally joining in one or the other.

The brothers talked of their own concerns and pursuits, but principally of those of the elder, whose temper was by much the most communicative, and who was always the greater talker. As a magistrate, he had generally some point of law to consult John about, or, at least, some curious anecdote to give; and as a farmer, as keeping in hand the home-farm at Donwell, he had to tell what every field was to bear next year, and to give all such local information as could not fail of being interesting to a brother whose home it had equally been the longest part of his life, and whose attachments were strong. The plan of a drain, the change of a fence, the felling of a tree, and the destination of every acre for wheat, turnips, or spring corn, was entered into with as much equality of interest by John, as his cooler manners rendered possible; and if his willing brother ever left him any thing to inquire about, his inquiries even approached a tone of eagerness.

While they were thus comfortably occupied, Mr Woodhouse was enjoying a full flow of happy regrets and fearful affection with his daughter.

<My poor dear Isabella,> said he, fondly taking her hand, and interrupting, for a few moments, her busy labours for some one of her five children—<How long it is, how terribly long since you were here! And how tired you must be after your journey! You must go to bed early, my dear—and I recommend a little gruel to you before you go.—You and I will have a nice basin of gruel together. My dear Emma, suppose we all have a little gruel.>

Emma could not suppose any such thing, knowing as she did, that both the Mr Knightleys were as unpersuadable on that article as herself;—and two basins only were ordered. After a little more discourse in praise of gruel, with some wondering at its not being taken every evening by every body, he proceeded to say, with an air of grave reflection,

<It was an awkward business, my dear, your spending the autumn at South End instead of coming here. I never had much opinion of the sea air.>

<Mr Wingfield most strenuously recommended it, sir—or we should not have gone. He recommended it for all the children, but particularly for the weakness in little Bella's throat,—both sea air and bathing.>

<Ah! my dear, but Perry had many doubts about the sea doing her any good; and as to myself, I have been long perfectly convinced, though perhaps I never told you so before, that the sea is very rarely of use to any body. I am sure it almost killed me once.>

<Come, come,> cried Emma, feeling this to be an unsafe subject, <I must beg you not to talk of the sea. It makes me envious and miserable;—I who have never seen it! South End is prohibited, if you please. My dear Isabella, I have not heard you make one inquiry about Mr Perry yet; and he never forgets you.>

<Oh! good Mr Perry—how is he, sir?>

<Why, pretty well; but not quite well. Poor Perry is bilious, and he has not time to take care of himself—he tells me he has not time to take care of himself—which is very sad—but he is always wanted all round the country. I suppose there is not a man in such practice anywhere. But then there is not so clever a man any where.>

<And Mrs Perry and the children, how are they? do the children grow? I have a great regard for Mr Perry. I hope he will be calling soon. He will be so pleased to see my little ones.>

<I hope he will be here to-morrow, for I have a question or two to ask him about myself of some consequence. And, my dear, whenever he comes, you had better let him look at little Bella's throat.>

<Oh! my dear sir, her throat is so much better that I have hardly any uneasiness about it. Either bathing has been of the greatest service to her, or else it is to be attributed to an excellent embrocation of Mr Wingfield's, which we have been applying at times ever since August.>

<It is not very likely, my dear, that bathing should have been of use to her—and if I had known you were wanting an embrocation, I would have spoken to—>

<You seem to me to have forgotten Mrs and Miss Bates,> said Emma, <I have not heard one inquiry after them.>

<Oh! the good Bateses—I am quite ashamed of myself—but you mention them in most of your letters. I hope they are quite well. Good old Mrs Bates—I will call upon her to-morrow, and take my children.—They are always so pleased to see my children.—And that excellent Miss Bates!—such thorough worthy people!—How are they, sir?>

<Why, pretty well, my dear, upon the whole. But poor Mrs Bates had a bad cold about a month ago.>

<How sorry I am! But colds were never so prevalent as they have been this autumn. Mr Wingfield told me that he has never known them more general or heavy—except when it has been quite an influenza.>

<That has been a good deal the case, my dear; but not to the degree you mention. Perry says that colds have been very general, but not so heavy as he has very often known them in November. Perry does not call it altogether a sickly season.>

<No, I do not know that Mr Wingfield considers it very sickly except\longdash>

<Ah! my poor dear child, the truth is, that in London it is always a sickly season. Nobody is healthy in London, nobody can be. It is a dreadful thing to have you forced to live there! so far off!—and the air so bad!>

<No, indeed—we are not at all in a bad air. Our part of London is very superior to most others!—You must not confound us with London in general, my dear sir. The neighbourhood of Brunswick Square is very different from almost all the rest. We are so very airy! I should be unwilling, I own, to live in any other part of the town;—there is hardly any other that I could be satisfied to have my children in: but we are so remarkably airy!—Mr Wingfield thinks the vicinity of Brunswick Square decidedly the most favourable as to air.>

<Ah! my dear, it is not like Hartfield. You make the best of it—but after you have been a week at Hartfield, you are all of you different creatures; you do not look like the same. Now I cannot say, that I think you are any of you looking well at present.>

<I am sorry to hear you say so, sir; but I assure you, excepting those little nervous head-aches and palpitations which I am never entirely free from anywhere, I am quite well myself; and if the children were rather pale before they went to bed, it was only because they were a little more tired than usual, from their journey and the happiness of coming. I hope you will think better of their looks to-morrow; for I assure you Mr Wingfield told me, that he did not believe he had ever sent us off altogether, in such good case. I trust, at least, that you do not think Mr Knightley looking ill,> turning her eyes with affectionate anxiety towards her husband.

<Middling, my dear; I cannot compliment you. I think Mr John Knightley very far from looking well.>

<What is the matter, sir?—Did you speak to me?> cried Mr John Knightley, hearing his own name.

<I am sorry to find, my love, that my father does not think you looking well—but I hope it is only from being a little fatigued. I could have wished, however, as you know, that you had seen Mr Wingfield before you left home.>

<My dear Isabella,>—exclaimed he hastily—<pray do not concern yourself about my looks. Be satisfied with doctoring and coddling yourself and the children, and let me look as I chuse.>

<I did not thoroughly understand what you were telling your brother,> cried Emma, <about your friend Mr Graham's intending to have a bailiff from Scotland, to look after his new estate. What will it answer? Will not the old prejudice be too strong?>

And she talked in this way so long and successfully that, when forced to give her attention again to her father and sister, she had nothing worse to hear than Isabella's kind inquiry after Jane Fairfax; and Jane Fairfax, though no great favourite with her in general, she was at that moment very happy to assist in praising.

<That sweet, amiable Jane Fairfax!> said Mrs John Knightley.—<It is so long since I have seen her, except now and then for a moment accidentally in town! What happiness it must be to her good old grandmother and excellent aunt, when she comes to visit them! I always regret excessively on dear Emma's account that she cannot be more at Highbury; but now their daughter is married, I suppose Colonel and Mrs Campbell will not be able to part with her at all. She would be such a delightful companion for Emma.>

Mr Woodhouse agreed to it all, but added,

<Our little friend Harriet Smith, however, is just such another pretty kind of young person. You will like Harriet. Emma could not have a better companion than Harriet.>

<I am most happy to hear it—but only Jane Fairfax one knows to be so very accomplished and superior!—and exactly Emma's age.>

\begin{figure}[tbph]
\centering
\includegraphics[width=.9\linewidth]{12tolerable}
\caption{She had never been able to get any thing tolerable}
\end{figure}

This topic was discussed very happily, and others succeeded of similar moment, and passed away with similar harmony; but the evening did not close without a little return of agitation. The gruel came and supplied a great deal to be said—much praise and many comments—undoubting decision of its wholesomeness for every constitution, and pretty severe Philippics upon the many houses where it was never met with tolerably;—but, unfortunately, among the failures which the daughter had to instance, the most recent, and therefore most prominent, was in her own cook at South End, a young woman hired for the time, who never had been able to understand what she meant by a basin of nice smooth gruel, thin, but not too thin. Often as she had wished for and ordered it, she had never been able to get any thing tolerable. Here was a dangerous opening.

<Ah!> said Mr Woodhouse, shaking his head and fixing his eyes on her with tender concern.—The ejaculation in Emma's ear expressed, <Ah! there is no end of the sad consequences of your going to South End. It does not bear talking of.> And for a little while she hoped he would not talk of it, and that a silent rumination might suffice to restore him to the relish of his own smooth gruel. After an interval of some minutes, however, he began with,

<I shall always be very sorry that you went to the sea this autumn, instead of coming here.>

<But why should you be sorry, sir?—I assure you, it did the children a great deal of good.>

<And, moreover, if you must go to the sea, it had better not have been to South End. South End is an unhealthy place. Perry was surprized to hear you had fixed upon South End.>

<I know there is such an idea with many people, but indeed it is quite a mistake, sir.—We all had our health perfectly well there, never found the least inconvenience from the mud; and Mr Wingfield says it is entirely a mistake to suppose the place unhealthy; and I am sure he may be depended on, for he thoroughly understands the nature of the air, and his own brother and family have been there repeatedly.>

<You should have gone to Cromer, my dear, if you went anywhere.—Perry was a week at Cromer once, and he holds it to be the best of all the sea-bathing places. A fine open sea, he says, and very pure air. And, by what I understand, you might have had lodgings there quite away from the sea—a quarter of a mile off—very comfortable. You should have consulted Perry.>

<But, my dear sir, the difference of the journey;—only consider how great it would have been.—An hundred miles, perhaps, instead of forty.>

<Ah! my dear, as Perry says, where health is at stake, nothing else should be considered; and if one is to travel, there is not much to chuse between forty miles and an hundred.—Better not move at all, better stay in London altogether than travel forty miles to get into a worse air. This is just what Perry said. It seemed to him a very ill-judged measure.>

Emma's attempts to stop her father had been vain; and when he had reached such a point as this, she could not wonder at her brother-in-law's breaking out.

<Mr Perry,> said he, in a voice of very strong displeasure, <would do as well to keep his opinion till it is asked for. Why does he make it any business of his, to wonder at what I do?—at my taking my family to one part of the coast or another?—I may be allowed, I hope, the use of my judgment as well as Mr Perry.—I want his directions no more than his drugs.> He paused—and growing cooler in a moment, added, with only sarcastic dryness, <If Mr Perry can tell me how to convey a wife and five children a distance of an hundred and thirty miles with no greater expense or inconvenience than a distance of forty, I should be as willing to prefer Cromer to South End as he could himself.>

<True, true,> cried Mr Knightley, with most ready interposition—<very true. That's a consideration indeed.—But John, as to what I was telling you of my idea of moving the path to Langham, of turning it more to the right that it may not cut through the home meadows, I cannot conceive any difficulty. I should not attempt it, if it were to be the means of inconvenience to the Highbury people, but if you call to mind exactly the present line of the path.... The only way of proving it, however, will be to turn to our maps. I shall see you at the Abbey to-morrow morning I hope, and then we will look them over, and you shall give me your opinion.>

Mr Woodhouse was rather agitated by such harsh reflections on his friend Perry, to whom he had, in fact, though unconsciously, been attributing many of his own feelings and expressions;—but the soothing attentions of his daughters gradually removed the present evil, and the immediate alertness of one brother, and better recollections of the other, prevented any renewal of it.